\section{\texorpdfstring{Semi-discrete and $\Wass_1$}{Semi-discrete and W1}}
\label{sec-semidiscr-w1}
\subsection{Semi-dual}\label{sec-semi-dual}
\paragraph{General measure semi-dual.}
Denote by $\Ee_0(f,g)$ the full dual objective $\int f\d\al+\int g\d\be$, extended by $-\infty$ outside $\Potentials(\c)$. For fixed $g$, dual feasibility is equivalent to $f\leq g^{\bar c}$, so positivity of $\al$ gives the unconstrained concave problem
\eql{\label{eq-semi-dual}
\MK_\c(\al,\be)
=
\usup{g \in \Cc(\Yy)}\Ee(g),
\qquad
\Ee(g)\eqdef\Ee_0(g^{\bar c},g)
=\int_\Xx g^{\bar c}\d\al+\int_\Yy g\d\be.
}
It is invariant under $g\mapsto g+s$.

\paragraph{Discrete semi-dual.}
For $\al=\sum_i\a_i\de_{x_i}$, $\be=\sum_j\b_j\de_{y_j}$, $\C_{ij}=c(x_i,y_j)$, and common mass $M$, eliminating the source potential gives
\eql{\label{eq-discrete-semi-dual}
\MKD_\C(\a,\b)
=
\umax{\gD\in\RR^m}\Ee(\gD),
\qquad
\Ee(\gD)
=
\Ee_0(\gD^{\bar\C},\gD)
=
\sum_i\a_i(\gD^{\bar\C})_i+\sum_j\b_j\gD_j,
\quad
(\gD^{\bar\C})_i=\min_j(\C_{ij}-\gD_j).
}
The objective is concave, piecewise affine, and gauge invariant. If $\sigma_\gD(i)\in\argmin_j(\C_{ij}-\gD_j)$, then $\b-\widehat\b(\gD)$ is a supergradient, where $\widehat\b_j(\gD)=\sum_{i:\sigma_\gD(i)=j}\a_i$.

\subsection{Auction Algorithm}\label{sec-auction-dual-ascent}
For the square uniform assignment problem, auction is derived from coordinate maximization of the discrete semi-dual, but uses bidder-specific $\epsilon$-overshoots to avoid jamming at nonsmooth ties. With $\gD$ denoting the target Kantorovich potential, set
\[
(\gD^{\bar\C})_i=\min_j(\C_{ij}-\gD_j),
\qquad
\Ee(\gD)=\frac1n\sum_i(\gD^{\bar\C})_i+\frac1n\sum_j\gD_j.
\]
The discrete Laguerre cells are
\[
\Laguerre_j^{\rm D}(\gD)=\{i:\C_{ij}-\gD_j\leq\C_{ik}-\gD_k\ \forall k\}.
\]
They are the restriction of Definition~\ref{def-laguerre-power-cells}. At points with unique minimizers,
\(
\partial_{\gD_j}\Ee(\gD)=(1-|\Laguerre_j^{\rm D}(\gD)|)/n.
\)

\paragraph{Bids and certificate.}
For an unassigned row $i$, choose its best and second-best targets $j_0,j_1$ and update
\eql{\label{eq-auction-dual-update}
\gD_{j_0}\leftarrow\gD_{j_0}
-\left[(\C_{i,j_1}-\gD_{j_1})-(\C_{i,j_0}-\gD_{j_0})+\epsilon\right].
}
Assign row $i$ to $j_0$ and unassign the previous owner of $j_0$, if any. The overshoot need not increase the nonsmooth semi-dual; its invariant is instead
\[
\P_{ij}=1\quad\Longrightarrow\quad
\C_{ij}-\gD_j\leq(\gD^{\bar\C})_i+\epsilon.
\]
If $R_\C=\max\C_{ij}-\min\C_{ij}$ and $\gD=0$ initially, the method terminates after at most $n(\lfloor R_\C/\epsilon\rfloor+1)$ bids and
\[
0\leq\frac1n\dotp{\C}{\P}-\min_{\P'\in\mathcal P_n^{\rm perm}}\frac1n\dotp{\C}{\P'}\leq\epsilon.
\]
The second inequality follows by applying Proposition~\ref{prop-assignment-dual-certificate} to the feasible pair $(\gD^{\bar\C},\gD)$ defined by the final target potential.
Dense implementation costs $O(n^2(1+R_\C/\epsilon))$. If $\C$ is integral, $\epsilon<1/n$ gives an exact assignment. Warm-starting auctions while halving $\epsilon$ yields an $\eta$-optimal assignment in $O(n^3(1+\log_+(R_\C/\eta)))$ operations.
\subsection{Semi-discrete}
\paragraph{Discrete target and Laguerre cells.}
One can adapt Definition~\ref{def-c-transform} of the $\bar c$-transform to this setting by restricting the minimization to the finite support $\{y_j\}_{j=1}^m$ of $\be$. Equivalently, one applies the same definition with the discrete target space $\Y=\{y_j\}_{j=1}^m$ and identifies a vector $\gD\in\RR^m$ with the function $g:\Y\to\RR$ defined by $g(y_j)=\gD_j$:

\eql{\label{eq-disc-c-transfo}
\foralls \gD \in \RR^m, \;
\foralls x \in \Xx, \quad
\gD^{\bar \c}(x) \eqdef \umin{j \in \range{m}} \c(x,y_j) - \gD_j.
}
\eql{\label{eq-semi-dual-discr}
\MK_\c(\al,\be) =
\umax{\gD \in \RR^m}
\Ee(\gD) \eqdef \Ee_0(\gD^{\bar c},\gD) =
\int_\Xx \gD^{\bar \c}(x) \d\al(x) + \sum_j \gD_j \b_j.
}
The objective satisfies $\Ee(\gD+s\ones)=\Ee(\gD)$ for every $s\in\RR$. One may therefore impose the gauge $\sum_j\gD_j=0$ without changing the problem.

\begin{defn}[Laguerre cells and power diagrams]\label{def-laguerre-power-cells}
For sites $(y_j)_{j=1}^m$ and weights $\gD\in\RR^m$, the Laguerre cell associated with $y_j$ is
\eql{\label{eq-laguerre-cells}
\Laguerre_{j}(\gD) \eqdef \enscond{x \in \Xx}{ \foralls j' \neq j, \c(x,y_j) - \gD_j \leq \c(x,y_{j'}) - \gD_{j'} }.
}
The cells cover $\Xx$; after arbitrary tie-breaking on common boundaries, they induce a disjoint partition. When $c(x,y)=\norm{x-y}^2$, this Laguerre decomposition is also called a power diagram. If $\gD$ is constant, it reduces to the ordinary Voronoi diagram.
\end{defn}
\paragraph{Mass balance.}
\eql{\label{eq-semi-disc-energy}
\Ee(\gD) = \sum_{j=1}^m \int_{\Laguerre_{j}(\gD)} \pa{ c(x,y_j) - \gD_j } \d\al(x) + \dotp{\gD}{\b}.
}
\begin{prop}[Gradient of the semi-discrete dual]\label{prop-semidiscrete-dual-gradient}
If the minimizing index in~\eqref{eq-disc-c-transfo} is unique for $\al$-almost every $x$, then $\Ee$ is differentiable at $\gD$ and
\(\foralls j \in \range{m}, \quad \nabla\Ee(\gD)_j = \b_j - \int_{\Laguerre_{j}(\gD)} \d\al.\)
\end{prop}
\begin{proof}
For $\al$-almost every $x$, the minimizing index in $\min_j c(x,y_j)-\gD_j$ is unique. If this index is $j(x)$, then the directional derivative in a direction $h\in\RR^m$ is
\[
\left.\frac{\d}{\d t}\right|_{t=0}
\min_j \bigl(c(x,y_j)-(\gD_j+t h_j)\bigr)
=
-h_{j(x)}.
\]
The corresponding difference quotients are bounded by $\norm{h}_\infty$. Dominated convergence therefore gives \(\d\Ee(\gD)[h] = -\sum_j h_j\int_{\Laguerre_j(\gD)}\d\al +\sum_j h_j\b_j,\)
which is the announced gradient formula.
\end{proof}
\paragraph{Stochastic optimization.}
\eql{\label{eq-semi-disc-energy-entropy}
\Ee(\gD)
=
\int_\Xx E(\gD,x)\d\al(x)
=
\EE_X(E(\gD,X)),
\qquad
E(\gD,x)\eqdef \gD^{\bar c}(x)+\dotp{\gD}{\b}.
}
\[
\nabla_{\gD}E(\gD,x)
=
\left(\b_j-\ones_{\Laguerre_j(\gD)}(x)\right)_{j=1}^m,
\]
\eql{\label{eq-sgd}
\gD^{(\ell+1)}
\eqdef
\gD^{(\ell)}+\tau_\ell\nabla_{\gD}E(\gD^{(\ell)},x_\ell).
}
\(j_\ell\in\argmin_j\bigl(c(x_\ell,y_j)-\gD_j^{(\ell)}\bigr),\)
\(\gD_j^{(\ell+1)} = \gD_j^{(\ell)} + \tau_\ell\bigl(\b_j-\ones_{\{j=j_\ell\}}\bigr).\)
\eql{\label{eq-step-size-sgd}
\sum_{\ell=0}^{+\infty}\tau_\ell=+\infty,
\qquad
\sum_{\ell=0}^{+\infty}\tau_\ell^2<+\infty,
}
\begin{prop}[Averaged stochastic semi-dual rate]\label{prop-semidiscrete-sgd-rate}
Let $\gD^\star$ maximize $\Ee$, set $R=\norm{\gD^{(0)}-\gD^\star}_2$, and suppose the stochastic supergradients are conditionally unbiased and satisfy $\norm{\nabla_\gD E(\gD,x)}_2\leq G$. For a horizon $L\geq1$, use the constant step $\tau=R/(G\sqrt L)$ and the average $\bar\gD_L=L^{-1}\sum_{\ell=0}^{L-1}\gD^{(\ell)}$. If $R=0$, the conclusion is immediate; otherwise,
\(\Ee(\gD^\star)-\EE\bigl[\Ee(\bar\gD_L)\bigr] \leq \frac{RG}{\sqrt L}.\)
\end{prop}
\begin{proof}
Conditionally on $\gD^{(\ell)}$, the stochastic supergradient is unbiased. Concavity and the squared-distance recursion give
\[
2\tau\,\EE\!\left[\Ee(\gD^\star)-\Ee(\gD^{(\ell)})\right]
\leq
\EE\norm{\gD^{(\ell)}-\gD^\star}_2^2
-
\EE\norm{\gD^{(\ell+1)}-\gD^\star}_2^2
+
\tau^2G^2.
\]
Summing from $0$ to $L-1$, discarding the final nonnegative squared distance, and using concavity at $\bar\gD_L$ yields
\[
\Ee(\gD^\star)-\EE\!\left[\Ee(\bar\gD_L)\right]
\leq
\frac{R^2}{2\tau L}+\frac{\tau G^2}{2}.
\]
The stated choice of $\tau$ gives the result.
\end{proof}
\subsection{Optimal Quantization}
\label{sec-optimal-quantization}
\paragraph{Free masses and prescribed weights.}
\eql{\label{eq-optimal-quantization}
\Qq_m(\al)
\eqdef
\umin{Y=(y_j)_{j=1}^m,\ \b\in\simplex_m}
\Wass_p\left(\al,\sum_{j=1}^m \b_j\de_{y_j}\right).
}
\begin{prop}[Quantization rate and curse of dimensionality]\label{prop-quantization-rate}
Let $\al$ be supported in a bounded subset $\Omega\subset\RR^d$ and assume that $\al=\rho\,\d x$ with $\rho\leq\rho_+<+\infty$. Then, for fixed $p\geq1$, there exist constants $0<c\leq C<+\infty$ such that
\(c\,m^{-1/d} \leq \Qq_m(\al) \leq C\,m^{-1/d}.\)
\end{prop}
\begin{proof}
Enclose $\Omega$ in a cube. With $k=\lfloor m^{1/d}\rfloor$, subdivide this cube into $k^d\leq m$ congruent subcubes and place one codepoint in each subcube meeting $\Omega$. Every source point then lies within distance $C/k\leq C'm^{-1/d}$ of a codepoint. The nearest-codepoint coupling proves the upper bound; unused codepoints may be assigned zero mass.

For the lower bound, fix any set $Y$ of at most $m$ codepoints and write $d_Y(x)=\min_j\norm{x-y_j}$. If $\omega_d$ is the volume of the Euclidean unit ball, then
\(
\al(\{d_Y\leq t\})\leq \rho_+m\omega_dt^d.
\)
Set $t_0=(2\rho_+m\omega_d)^{-1/d}$. Then $\al(\{d_Y>t\})\geq1/2$ for $0<t<t_0$, and the layer-cake formula gives \(\int d_Y(x)^p\d\al(x) = \int_0^{+\infty} p t^{p-1}\al(\{d_Y>t\})\d t \geq \frac{t_0^p}{2} \simeq c m^{-p/d}.\)
Taking the $p$-th root and minimizing over $Y$ proves the lower bound.
\end{proof}
\paragraph{Lloyd algorithm.}
\begin{prop}[Free masses give Voronoi cells]\label{prop-free-masses-voronoi}
For the cost $c(x,y)=d(x,y)^p$, fix distinct codepoints $Y=(y_j)_{j=1}^m$. Duplicate codepoints can be merged beforehand. Minimizing over the weights $\b\in\simplex_m$ gives
\[
\min_{\b\in\simplex_m}
\Wass_p^p\left(\al,\sum_j \b_j\de_{y_j}\right)
=
\int_\X \min_{1\leq j\leq m} c(x,y_j)\d\al(x).
\]
An optimal coupling is induced by sending each $x$ to a nearest codepoint. The corresponding cells are the Voronoi cells
\(\VV_j(Y) \eqdef \enscond{x}{\foralls j',\ c(x,y_j)\leq c(x,y_{j'})},\)
up to arbitrary tie-breaking on common boundaries.
\end{prop}
\begin{proof}
For any coupling between $\al$ and a measure supported on $Y$, the conditional destination of a point $x$ belongs to $Y$, hence its conditional cost is at least $\min_j c(x,y_j)$. Integrating gives the lower bound. Conversely, choose a measurable nearest-codepoint map $T_Y(x)\in\argmin_j c(x,y_j)$, breaking ties measurably, and set $\b_j=\al(T_Y^{-1}(y_j))$. Then $(T_Y)_\sharp\al=\sum_j \b_j\de_{y_j}$ and the induced transport reaches the displayed lower bound.
\end{proof}
\(\Qq_m(\al)^p = \min_Y \Ff(Y), \qquad \Ff(Y) \eqdef \int_\X \min_{1\leq j\leq m} c(x,y_j)\d\al(x).\)
\(y_j \in \uargmin{y} \int_{\VV_j(Y)} c(x,y)\d\al(x).\)
\(y_j = \frac{\int_{\VV_j(Y)} x\d\al(x)} {\int_{\VV_j(Y)} \d\al}.\)
\paragraph{Continuous Lloyd flow.}
\(a_j(Y)=\al(\VV_j(Y)), \qquad b_j(Y)=\frac{1}{a_j(Y)}\int_{\VV_j(Y)}x\,\d\al(x)\)
\(y_j^{(k+1)}=y_j^{(k)}+\tau\big(b_j(Y^{(k)})-y_j^{(k)}\big), \qquad 0<\tau\leq 1,\)
\(\dot y_j(t)=b_j(Y_t)-y_j(t).\)
\(\nabla_{y_j}\Ff(Y) = 2\int_{\VV_j(Y)}(y_j-x)\d\al(x) = 2a_j(Y)(y_j-b_j(Y)),\)
\(\dot y_j(t) = -\frac{1}{2a_j(Y_t)}\nabla_{y_j}\Ff(Y_t).\)
\(\frac{\d}{\d t}\Ff(Y_t) = -2\sum_j a_j(Y_t)\norm{b_j(Y_t)-y_j(t)}^2 \leq 0.\)
If \(\eta_t=\sum_j w_j\de_{y_j(t)}\) carries fixed positive weights, independent of the Voronoi masses \(a_j(Y_t)\), Proposition~\ref{prop-lagrangian-flow-continuity} turns this labelled particle ODE into the weak continuity equation

\(\partial_t\eta_t+\diverg(v_t\eta_t)=0, \qquad v_t(y_j(t))=b_j(Y_t)-y_j(t),\)
in the sense of the measure evolutions of Section~\ref{sec-dynamic-optimal-transport}. The weights \(w_j\) in this transport equation are auxiliary weights for the moving labelled particles; they are not the Voronoi masses used to define the quantization energy.

\(\nu_{Y_t}=\sum_j a_j(Y_t)\de_{y_j(t)},\)
\(\partial_t\nu_{Y_t}+\diverg(v_t\nu_{Y_t}) = \sum_j \dot a_j(Y_t)\de_{y_j(t)}, \qquad v_t(y_j(t))=\dot y_j(t).\)
\paragraph{Mean-field limit and ultrafast diffusion.}
There is nevertheless a precise high-resolution continuum theory when the number \(m\) of codepoints tends to infinity. If \(\al=\rho\d x\), the limiting Eulerian variable is the density \(\sigma\) of sites, meaning heuristically that \(m\sigma(x)\d x\) codepoints lie in \(\d x\).

\(\Gg_\rho(\sigma) \eqdef \int_\Omega \rho(x) \sigma(x)^{-r}\d x, \qquad \int_\Omega \sigma\,\d x=1,\)
\begin{align*}
m^r\Ff_p(Y^{(m)})
&\simeq C_{p,d}\Gg_\rho(\sigma),\\
m^r\Qq_m(\al)^p
&\longrightarrow C_{p,d}\min_\sigma\Gg_\rho(\sigma).
\end{align*}
\[
\partial_t \sigma
=
-r\,\diverg\!\left(
\sigma\nabla\!\left(\frac{\rho}{\sigma^{r+1}}\right)
\right),
\]
\[
\partial_t u
=
-\frac{r+1}{\omega}\diverg\!\left(\omega\nabla(u^{-r})\right),
\]
\paragraph{Quantization with fixed equal weights.}
\(\nu_Y\eqdef \frac1m\sum_{j=1}^m\de_{y_j}, \qquad \Ff_{\rm eq}(Y) \eqdef \frac12\Wass_2^2(\al,\nu_Y),\)
\(\bar x_j(Y) \eqdef m\int_{C_j(Y)} x\d\al(x).\)
\(\nabla_{y_j}\Ff_{\rm eq}(Y) = \frac1m\bigl(y_j-\bar x_j(Y)\bigr).\)
As long as the optimal matching between two nearby labelled configurations preserves labels, the \(\Wass_2\) metric on equal-weight empirical measures induces the local particle metric \(g_Y(U,V)=m^{-1}\sum_j\dotp{u_j}{v_j}\). Hence the associated \(\Wass_2\) gradient flow is the coupled system

\(\dot y_j(t) = \bar x_j(Y_t)-y_j(t), \qquad j=1,\ldots,m.\)
Equivalently, \(\nu_{Y_t}\) satisfies a continuity equation with velocity \(v_t(y_j(t))=\bar x_j(Y_t)-y_j(t)\). This is the so-called finite-particle \(\Wass_2\) gradient-flow viewpoint developed more systematically in Chapter~\ref{sec-wasserstein-gradient-flows}; the fixed-weight cells are Laguerre cells rather than the free-mass Voronoi cells used by Lloyd's method.

\paragraph{Equal-weight quantization on the line.}
\begin{prop}[One-dimensional equal-weight quantization]\label{prop-1d-equal-weight-quantization}
Let $\al\in\Mm_+^1(\RR)$ have finite second moment and quantile function $Q=\cumul{\al}^{-1}$. For
\[
\Qq^{\mathrm{eq}}_m(\al)^2
\eqdef
\min_{y_1\leq\cdots\leq y_m}
\Wass_2^2\left(\al,\frac1m\sum_{i=1}^m\de_{y_i}\right),
\]
set $I_i=((i-1)/m,i/m]$. Then the sorted minimizer is unique and its $i$th atom is \(y_i^\star = m\int_{I_i} Q(u)\d u,\)
and
\(\Qq^{\mathrm{eq}}_m(\al)^2 = \sum_{i=1}^m \int_{I_i} \abs{Q(u)-y_i^\star}^2\d u.\)
If $Q\in C^1([0,1])$, then
\(m^2\Qq^{\mathrm{eq}}_m(\al)^2 \longrightarrow \frac1{12}\int_0^1\abs{Q'(u)}^2\d u.\)
\end{prop}
\begin{proof}
After sorting the atoms, the quantile formula for $\Wass_2$ gives
\[
\Wass_2^2\left(\al,\frac1m\sum_{i=1}^m\de_{y_i}\right)
=
\sum_{i=1}^m
\int_{I_i}\abs{Q(u)-y_i}^2\d u.
\]
The minimization therefore decouples over the intervals $I_i$, and the best constant approximation of $Q$ on $I_i$ in $L^2$ is its average on $I_i$. Since $Q$ is nondecreasing, these interval averages are nondecreasing, so they satisfy the sorting constraint. Strict convexity in each scalar $y_i$ proves uniqueness and gives the displayed energy.

Denote by $\bar Q_{I_i}=m\int_{I_i}Q(u)\d u$ the interval average. If $Q$ is $C^1$, set $h=1/m$ and write $I_i=(a_i,a_i+h]$. Uniform Taylor expansion gives, for $v\in[0,1]$,
\(Q(a_i+hv)=Q(a_i)+hvQ'(a_i)+h\,r_i(v), \qquad \max_i\sup_{v\in[0,1]}\abs{r_i(v)}\to0.\)
Subtracting the average over $v\in[0,1]$ and integrating gives \(\int_{I_i}\abs{Q(u)-\bar Q_{I_i}}^2\d u = \frac{h^3}{12}\abs{Q'(a_i)}^2+o(h^3),\)
uniformly in $i$. Summing over $i$ yields a Riemann sum for $\frac1{12}\int_0^1\abs{Q'(u)}^2\d u$.
\end{proof}
Thus the common deterministic rule $m^{-1}\sum_i\de_{Q((i-1/2)/m)}$ should be read as a midpoint approximation of the optimal bin-average formula. It has the same leading squared error under the same smoothness assumptions.

\(\int_{I_i}|Q(u)-z|^2\d u = \int_{I_i}|Q(u)-y_i^\star|^2\d u + \frac1m|z-y_i^\star|^2,\)
\begin{prop}[Quantile-process asymptotics for random placement]\label{prop-1d-random-quantile-process}
Let $\al\in\Mm_+^1(\RR)$ have quantile $Q\in C^1([0,1])$, and let $\widehat\al_m=m^{-1}\sum_{i=1}^m\de_{X_i}$ with $X_i$ i.i.d. with law $\al$. If $B$ denotes the standard Brownian bridge on $[0,1]$, then
\(m\,\Wass_2^2(\al,\widehat\al_m) \overset{\mathrm{law}}{\longrightarrow} \int_0^1 B(u)^2\abs{Q'(u)}^2\d u,\)
and, in expectation,
\[
m\,\EE\!\left[\Wass_2^2(\al,\widehat\al_m)\right]
\longrightarrow
\int_0^1 u(1-u)\abs{Q'(u)}^2\d u.
\]
\end{prop}
\begin{proof}
Let $\widehat Q_m$ be the empirical quantile function. The one-dimensional formula gives \(\Wass_2^2(\al,\widehat\al_m) = \int_0^1\abs{\widehat Q_m(u)-Q(u)}^2\d u.\)
Write $X_i=Q(U_i)$ with $U_i$ i.i.d. uniform on $[0,1]$, and let $\widehat U_m^{-1}$ be the empirical quantile function of the $U_i$. Then $\widehat Q_m=Q\circ \widehat U_m^{-1}$. The classical uniform quantile-process theorem gives
\(\sqrt m\,(\widehat U_m^{-1}-\Id) \overset{\mathrm{law}}{\longrightarrow} -B \quad\text{in }L^2(0,1).\)
Because $Q'$ is uniformly continuous and $\widehat U_m^{-1}\to\Id$ uniformly in probability, the functional delta method gives \(\sqrt m\,(\widehat Q_m-Q) \overset{\mathrm{law}}{\longrightarrow} -B\,Q' \quad\text{in }L^2(0,1).\)
Since $B$ and $-B$ have the same law, the sign disappears in the squared $L^2$ norm. The continuous mapping theorem gives the distributional convergence. Standard fourth-moment bounds for the uniform quantile process, together with the boundedness of $Q'$, give uniform integrability of the squared $L^2$ norms and hence convergence of expectations. Since $\EE[B(u)^2]=u(1-u)$, Fubini's theorem gives the displayed expectation limit.
\end{proof}
\subsection{Wasserstein-1 norm}
\label{sec-W1}
\paragraph{\texorpdfstring{$c$-transform for $\Wass_1$.}{c-transform for W1.}}
Assume that $d$ is a distance on $\X=\Y$ and take the ground cost $c(x,y)=d(x,y)$. We denote the Lipschitz constant of $f\in\Cc(\X)$ by

\begin{defn}[Lipschitz constant]\label{def-lipschitz-constant}
For a function $f:\X\to\RR$ on a metric space $(\X,d)$, its Lipschitz constant is
\eql{\label{eq-lip-constant}
\Lip(f)
\eqdef
\sup\enscond{\frac{|f(x)-f(y)|}{d(x,y)}}{x\neq y}.
}
The function is $1$-Lipschitz when $\Lip(f)\leq1$.
\end{defn}
\begin{prop}[$c$-transforms and $1$-Lipschitz functions]\label{prop-w1-c-transform-lipschitz}
Suppose $\X=\Y$ and $c(x,y)=d(x,y)$. Then there exists $g$ such that $f=g^c$ if and only if $\Lip(f)\leq1$. Furthermore, if $\Lip(f)\leq1$, then $f^c=-f$.
\end{prop}
\begin{proof}
First suppose $f=g^c$ for some $g$. For $x,y\in\X$,
\[
|f(x)-f(y)|
=
\left|\inf_z[d(x,z)-g(z)]-\inf_z[d(y,z)-g(z)]\right|
\leq
\sup_z |d(x,z)-d(y,z)|
\leq d(x,y),
\]
where the last inequality is the reverse triangle inequality. Thus $\Lip(f)\leq1$.

If $\Lip(f)\leq1$, then $f(x)\leq f(y)+d(x,y)$, so $d(x,y)-f(x)\geq -f(y)$ for all $x$ and hence $f^c(y)\geq -f(y)$. Taking $x=y$ gives $f^c(y)\leq -f(y)$. Therefore $f^c=-f$. Applying the same property to $-f$ gives $(-f)^c=f$, so every $1$-Lipschitz function is $c$-concave.
\end{proof}
\eql{\label{eq-w1-metric}
\Wass_1(\al,\be)
=
\umax{f}
\enscond{\int_\X f\d(\al-\be)}{\Lip(f)\leq1}
=:
\norm{\al-\be}_{W_1}.
}
For a discrete signed measure $\al-\be=\sum_{k=1}^N r_k\de_{z_k}$ on distinct sites, with $\sum_k r_k=0$,~\eqref{eq-w1-metric} becomes the finite-dimensional linear program

\eql{\label{eq-w1-discr}
\Wass_1(\al,\be)
=
\umax{(f_k)_k}
\enscond{\sum_k f_k r_k}{\foralls k,\ell,\ |f_k-f_\ell|\leq d(z_k,z_\ell)}.
}
This linear program can be solved by generic interior-point or first-order methods. If $N$ support points are involved, however, it still contains $O(N^2)$ Lipschitz constraints, mirroring the $O(nm)$ coupling variables of the original discrete Kantorovich LP; the dual formulation alone does not remove the all-pairs structure.

\(\Wass_1(\al,\be) = \umax{(f_k)_k} \enscond{\sum_k f_k r_k}{\foralls k,\ |f_{k+1}-f_k|\leq z_{k+1}-z_k}.\)
\paragraph{\texorpdfstring{$\Wass_1$ on Euclidean spaces.}{W1 on Euclidean spaces.}}
Let $\al,\be\in\Mm_+^1(\RR^d)$ have finite first moments and set $\xi=\al-\be$. Rademacher's theorem identifies $1$-Lipschitz functions with locally Sobolev functions whose Euclidean gradient has norm at most one almost everywhere. Hence

\eql{\label{eq-w1-cont}
\Wass_1(\al,\be) =
\usup{f\in W_{\mathrm{loc}}^{1,\infty}(\RR^d)}
\enscond{ \int_{\RR^d} f\d\xi }{ \norm{\nabla f}_{L^\infty(\RR^d)} \leq 1 }.
}
The correct flux space is generally not $L^1(\RR^d;\RR^d)$. Let $\Mm(\RR^d;\RR^d)$ denote the finite $\RR^d$-valued Radon measures, write $|\flow|$ for total variation, and define divergence distributionally by

\(\langle\diverg\flow,\varphi\rangle \eqdef -\int_{\RR^d}\dotp{\nabla\varphi}{\d\flow}, \qquad \varphi\in C_c^1(\RR^d).\)
\begin{prop}[Euclidean Beckmann formula]\label{prop-euclidean-beckmann}
For $\al,\be\in\Mm_+^1(\RR^d)$ with finite first moments,
\eql{\label{eq-w1-cont-div}
\Wass_1(\al,\be) =
\umin{\flow\in\Mm(\RR^d;\RR^d)}
\enscond{|\flow|(\RR^d)}{\diverg\flow=\al-\be}.
}
If $\flow=w\,\d x$ has a Lebesgue density, its cost reduces to $\int_{\RR^d}\norm{w(x)}_2\d x$.
\end{prop}
\begin{proof}
Let $\flow$ be feasible. For a smooth $1$-Lipschitz test function, \(\int f\d(\al-\be) = -\int\dotp{\nabla f}{\d\flow} \leq |\flow|(\RR^d).\)
Truncation followed by mollification approximates normalized Lipschitz functions by smooth test functions without increasing the gradient bound asymptotically. Equation~\eqref{eq-w1-cont} therefore gives $\Wass_1(\al,\be)\leq|\flow|(\RR^d)$.

Conversely, let $\pi$ be an optimal coupling and define a vector measure by its action on $\zeta\in C_c(\RR^d;\RR^d)$:
\(\int\dotp{\zeta(z)}{\d\flow(z)} \eqdef \int_{\RR^d\times\RR^d}\int_0^1 \dotp{\zeta((1-t)x+ty)}{y-x}\d t\d\pi(x,y).\)
For $\varphi\in C_c^1(\RR^d)$, the fundamental theorem of calculus gives \(\langle\diverg\flow,\varphi\rangle = -\int\!\int_0^1 \dotp{\nabla\varphi((1-t)x+ty)}{y-x}\d t\d\pi(x,y) = \int\varphi\d(\al-\be).\)
Moreover, $|\flow|(\RR^d)\leq\int\norm{x-y}_2\d\pi(x,y)=\Wass_1(\al,\be)$. Together with the first inequality, this proves equality and attainment.
\end{proof}
\paragraph{Graph distances and Beckmann flows.}
\begin{defn}[Graph geodesic distance]\label{def-graph-geodesic-distance}
Let $G=(V,E)$ be a connected finite graph with positive edge lengths $(\ell_e)_{e\in E}$. The graph geodesic distance between two vertices is \(d_G(i,j)=\min_{\gamma:i\leadsto j}\sum_{e\in\gamma}\ell_e.\)
The minimum is over all paths $\gamma$ joining $i$ to $j$.
\end{defn}
\begin{prop}[$\Wass_1$ and Beckmann flow on a graph]\label{prop-graph-w1-beckmann}
Let $G=(V,E)$ be a connected finite graph with positive edge lengths $(\ell_e)_{e\in E}$ and graph geodesic distance $d_G$.
For two probability vectors $\a,\b$ on $V$, set $r=\a-\b$ and orient each edge $e=(i,j)$. If
\((\nabla_G f)_e=f_j-f_i,\qquad \operatorname{div}_G=-\nabla_G^*\)
are the finite-difference gradient and its negative adjoint, then a positive flow on the oriented edge $i\to j$ has positive divergence at $i$ and negative divergence at $j$. With this convention,
\[
\Wass_{1,G}(\a,\b)
=
\max_f \enscond{\sum_{i\in V} f_i r_i}{|f_i-f_j|\leq \ell_e\quad\forall e=(i,j)}
=
\min_m \enscond{\sum_{e\in E}\ell_e |m_e|}{\operatorname{div}_G m=r}.
\]
The vector $m_e$ is an oriented edge flow, and the constraint $\operatorname{div}_G m=r$ is conservation of mass at each vertex.
\end{prop}
\begin{proof}
The edge constraints imply, by summing along every path and then minimizing over paths, that $|f_i-f_j|\leq d_G(i,j)$ for all vertices. Conversely, any $1$-Lipschitz function for $d_G$ satisfies $|f_i-f_j|\leq d_G(i,j)\leq\ell_e$ on every edge. The first equality is therefore the Kantorovich--Rubinstein formula on the metric space $(V,d_G)$.

For the second equality, write the graph Beckmann problem and dualize its equality constraint with a potential $f$: \(\inf_m \sum_e\ell_e|m_e| +\sup_f \sum_i f_i(r_i-(\operatorname{div}_G m)_i).\)
Using $\operatorname{div}_G=-\nabla_G^*$, the coupling term is $\sum_e m_e(\nabla_G f)_e$. The minimization over each scalar flow $m_e$ is finite exactly when $|(\nabla_G f)_e|\leq\ell_e$, and is then equal to zero. The dual problem is therefore precisely the graph Lipschitz dual above. Strong duality holds because this is a finite-dimensional linear program with a nonempty feasible set: connectedness and $\sum_i r_i=0$ allow the signed surplus to be routed along paths. This proves the graph Beckmann formula.
\end{proof}
